\chapter{The Field Closes}

\section{The second variation}

Part V opens on the field, and the field, on this face, is the book's own climax --- the place where the residue the physical
gauge has carried since Chapter 3 finally closes into the object physics calls electromagnetism. The chapter's first claim is
structural and exact: the field is the residue's \emph{second} variation. The first variation gave the equations of the last
chapter; the second variation --- the curvature of the action, $\delta^2 S$ --- is the field, and it closes into a definite,
finite shape the coupling will be asked of. This section builds the second variation and shows it is the field.

The mathematics is the Hessian of the action. The first variation vanishes on the realized configuration, $\delta S = 0$;
the second variation is what remains, the quadratic form
\[
  \delta^2 S \;=\; \int \delta\phi \;\frac{\delta^2 S}{\delta\phi\,\delta\phi}\; \delta\phi \;\mathrm{d}^4x ,
\]
whose kernel is the Jacobi operator --- the curvature of the space of configurations about the realized one. This is
curvature in the same sense as every earlier residue: it measures how the action bends as the configuration is varied, the
second-order failure of the extremal to be flat. And for the electromagnetic action the second variation \emph{is} the field
strength's own dynamics: varying $A$ twice returns the Maxwell operator, and the field $F = \mathrm{d}A$ is read off the
curvature the second variation exposes. The field is not a new object; it is the residue read at second order, the curvature
of the action about its extremal.

The construction renders this second variation not as a continuum but as a small, fixed number of \emph{channels}, and the
number is the construction's, not the book's: it enumerates exactly four and stops. Each channel is a ratio of the kind
Chapter 6 built --- a floored magnitude with an orientation --- so the field, at this tier, is four ratios: four finite
readings of the curvature, held as four remainders, with no continuum anywhere in them. The four channels are where the
electromagnetic coupling is assembled, the four contractions of the second variation that the construction carries. That the
field is four finite channels rather than a smooth quantity to be integrated is the load-bearing finiteness of the whole
chapter: the coupling, when asked, will be asked of these four channels at their floor, never of a continuous field.

It is worth marking this four against an earlier one. In Chapter 7 the residue was read four ways, and the book was careful
that \emph{that} four was its own --- the count of the gauge's projections, an expository choice. Here the four is different:
these channels are enumerated in the construction itself, four contractions it names and holds. So the expository four
faces of Part III close, at the field, onto four channels the construction actually carries --- a satisfying seam, but the distinction must be kept
exact. The faces were the book's four; the channels are the construction's four; that they meet is a genuine confirmation and
not a proof that the faces were derived. Two fours, meeting honestly at the field.

In the two verbs, the field is a fourfold funge of the residue's second variation, each channel a bag with a remainder ---
the residue tanged out of each exactly as the benches tanged it in the last part. The continuous field physics integrates is
a tange this face declines to complete; what it keeps is the four finite funges and the residue in each. So the field opens
the last part not as a smooth infinity but as four finite readings of the residue's curvature, the same leftover one order
deeper, held four ways and no way rounded. The next section names the field's other half --- the magnetic side that must be
present before the coupling can be asked --- and keeps it finite too.

\section{The magnetic side must be present}

The second variation gave the field its electric channels. But a field has two sides, and the coupling cannot be asked until
both are closed. This short section names the magnetic side, shows that the construction carries it and carries it finite,
and establishes that the field is whole --- both sides present, both floored --- before the last chapter asks the coupling.
The magnetic side is not an afterthought; it is the half of Maxwell's system that closes the field, and without it the
question of the coupling is not yet well-posed.

Maxwell's equations close exactly when both sides are present. Written in the language of forms, the two halves are
\[
  \mathrm{d}F = 0, \qquad \mathrm{d}\star F = \star J ,
\]
the first the magnetic side --- no monopoles, the field's flux conserved --- and the second the electric side, the source
equation. The magnetic side $\mathrm{d}F = 0$ is the Bianchi identity, the statement that $F$ is closed, that the residue's
loop face has no source of its own. It is exactly what makes $F = \mathrm{d}A$ possible, exactly the closure that lets the
holonomy of Chapter 3 be defined. Without the magnetic side the field does not close: there is a potential but no guarantee
its curvature is source-free, no well-defined loop to read. The magnetic side must be present because it is the closure
itself.

The construction carries the magnetic side as a Meissner-type reading, and carries it \emph{finite}, and this finiteness is
marked with care. The physical arrangement is the field-expulsion of Chapter 10 --- a screening that expels the applied field
over a penetration depth --- and the construction renders it not as a continuous London equation integrated over all space
but as a finite three-rung shape: three discrete stages of screening, not a smooth exponential to infinity. This is the same
discipline the whole book keeps, now applied to the field's magnetic half. The reader must not read the magnetic side as the
continuum Maxwell system physics integrates; it is a finite, three-rung ratio, an applied field over a shielding response
counted in discrete stages. The magnetic side is present, it closes the field, and it is finite --- marked as three rungs,
not a continuum.

With both sides present the field is whole, and the place they meet is the object the whole instrument was built to reach.
Where the electric channels and the magnetic side are set against each other --- where the two closed halves of the field
determine a single dimensionless ratio --- is the electromagnetic coupling. The construction has, at this point, everything
it needs to ask for it: four electric channels, a finite magnetic side, both floored, both in hand. And here the book does
what it has done since its first pages. It names the coupling --- the meeting of the two sides --- and carries not one digit
of its value. The field is closed; the coupling is named; its magnitude is withheld, held blind here as strictly as
everywhere, waiting one chapter more.

In the two verbs, the magnetic side funges into its finite three-rung ratio, the residue tanged out of it exactly as the
electric residue was. The coupling itself is a funge the instrument deliberately leaves open: the two sides bagged toward a
single number the chapter will not complete to a value. It is the same refusal the bracket has always been --- a bag held
without pinning, a number named without being fixed. So the field closes with both sides finite and in hand, the coupling
named and un-numbered between them --- everything assembled to ask, and nothing of the answer given. The last section
performs the asking, and delivers, as promised, the question and not the magnitude.

\section{Asking, not delivering}

The field is whole: both sides finite, both in hand, the coupling named between them without a number. This closing section
performs the one act the assembled field is for --- it \emph{asks} the coupling --- and then does the thing the whole book
has trained the reader to expect. It does not answer. The field's closure poses the question, fully and with everything in
place, and hands over the question, not the magnitude. This is the discipline at its most load-bearing, because it is the
chapter most tempted to deliver, and it refuses.

The asking is a contraction. The instrument gathers the four electric channels and the finite magnetic side and draws them
together toward a single dimensionless number --- the coupling --- by the contraction the construction carries for exactly
this purpose. In the field's own language the coupling is the ratio that the two closed sides determine between them, the
pure number left when the field's electric and magnetic halves are set against each other and all the dimensional scales
cancel. That gathering-and-contracting \emph{is} the asking. To assemble the field and contract it toward the coupling is to
put the question --- what is the coupling? --- in the only form the instrument can put it: an apparatus, built and aimed, a
contraction poised on the closed field.

The discipline that keeps the asking honest is the order of operations, and it is the anti-crank heart of the method. The
calibration behind the contraction is fixed \emph{before} the coupling is read, not fitted \emph{after} it. The instrument
does not build its contraction around a number it already wanted and then present that number as a discovery; it assembles
the field first, from the two finite sides, and only then contracts and asks what the assembly yields. That order is the
whole difference between an instrument and a fake. A fake tunes its apparatus to the answer it means to announce; this
instrument builds its apparatus before it looks and takes whatever the looking gives. Build, then ask --- never ask, then
build to match. The four channels and the three-rung side were closed with no target in the room, and the contraction reads
them with none.

And having asked, honestly, the closure does not deliver. The coupling's value is the next chapter's, held blind here as
strictly as on every page before. The field's closure poses the question --- everything assembled, the contraction aimed ---
and stops there, handing the question forward and keeping no number for itself. Nor will the next chapter, when it answers,
hand back a point: the answer will be a bracket, floored at the instrument's own resolution and returning to the same
interval every time it is asked. The magnitude the closure withholds is withheld not because the instrument is coy but
because the honest answer, when it comes, is not a magnitude at all. It is a bound. The field closes to ask; the bracket, one
chapter on, is what answers.

In the two verbs, the contraction funges the whole field into a single question --- the two sides, the four channels and the
three-rung ratio, bagged together as ``the coupling.'' And the instrument, having funged them into the question, declines the
last step: it will not tange the question down to a value. Asking is the honest funge; delivering a magnitude would be the
over-tange the book has refused since it first chose the bag over the pin. So the field's climax closes with the instrument
complete and the question posed --- the field contracted to the coupling, the coupling held open, and the number, one chapter
from here, still un-taken. Everything to ask with, and nothing yet answered.
